% =============================================================
% Chapitre 4 — Mise en œuvre (~13 pages)
% Contenu dérivé du LLD (conception cible détaillée) — À CONFRONTER
% à 06_ECARTS_IMPLEMENTATION.md avant validation : ce chapitre ne
% doit décrire, dans sa version finale, que ce qui est réellement
% déployé, pas seulement ce qui était prévu.
% =============================================================
\chapter{Mise en œuvre}
\label{ch:mise-en-oeuvre}

\begin{attentedoc}
Chapitre rédigé à partir du document de conception détaillée (LLD). Chaque sous-section
ci-dessous doit être confrontée à \texttt{06\_ECARTS\_IMPLEMENTATION.md} avant validation
finale : reformuler au passé (\og a été déployé \fg{}) uniquement ce qui est confirmé
réellement en place ; garder le conditionnel ou signaler explicitement un écart sinon.
\end{attentedoc}

\section{Stratégie d'environnements et provisionnement}

L'ensemble du socle --- des sept couches jusqu'au pipeline de livraison lui-même --- est
décrit exclusivement sous forme de code d'infrastructure, organisé en un module réutilisable
par couche. Aucune ressource n'est créée manuellement dans la console du fournisseur cloud :
toute création passe par un pipeline dédié qui formate, valide, planifie, scanne, fait
approuver puis applique les changements.

\begin{keybox}[title=Principe retenu]
Une image applicative est construite et analysée \textbf{une seule fois}, puis
\emph{promue} d'un environnement à l'autre par son identifiant de contenu (empreinte)
plutôt que reconstruite à chaque environnement. Ce qui s'exécute en production est, à
l'octet près, ce qui a été analysé en amont --- reconstruire annulerait la valeur des
analyses de sécurité.
\end{keybox}

\todo{Confirmer l'état réel de la stratégie de promotion (dev $\rightarrow$ recette
$\rightarrow$ production) au moment de la rédaction --- certains environnements peuvent ne
pas encore avoir été créés.}

\section{Déploiement de la couche identité et du réseau}

La couche identité et la couche réseau sont posées avant tout le reste, conformément à
l'ordre de construction retenu (\og de l'intérieur vers l'extérieur \fg{}) : les fondations
d'identité conditionnent la sécurité de toutes les couches suivantes. Cette couche comprend
la création des comptes de service, la définition des rôles associés, le stockage des
secrets applicatifs (un secret par valeur, jamais un secret dérivé d'un autre), et la
configuration du réseau privé avec un accès contrôlé aux ressources de données.

\section{Déploiement des workloads et du point d'entrée}

Les services applicatifs sont déployés sur une plateforme d'exécution serverless, avec un
point d'entrée unique filtré par un pare-feu applicatif. Chaque service ne peut être atteint
que via ce point d'entrée --- une tentative d'accès direct à l'adresse interne d'un service
est un des tests de validation prévus au chapitre~\ref{ch:tests}.

\section{Mise en œuvre du pipeline DevSecOps}

Le pipeline de livraison applique une séquence de portes bloquantes avant tout déploiement :

\begin{enumerate}
  \item authentification du pipeline par fédération d'identité (aucune clé stockée) ;
  \item détection de secrets accidentellement commités --- échec immédiat si un secret est
        détecté ;
  \item analyse statique du code ;
  \item construction de l'image, à partir d'une image de base identifiée par empreinte ;
  \item analyse de l'image et de ses dépendances (vulnérabilités connues) ;
  \item publication dans un registre unique de confiance ;
  \item déploiement référencé par empreinte de contenu, jamais par une étiquette mutable.
\end{enumerate}

Un second pipeline, distinct mais authentifié selon le même principe de fédération
d'identité, applique une séquence équivalente à l'infrastructure elle-même : formatage,
validation, plan de changement soumis à relecture humaine, analyse de sécurité de la
configuration d'infrastructure, puis application.

\section{Déploiement de la détection et de l'enrichissement sémantique}

Les journaux d'audit et applicatifs de l'ensemble des couches sont collectés puis filtrés
avant ingestion dans l'entrepôt de détection, afin de maîtriser le volume ingéré --- poste
de coût potentiellement dominant s'il n'est pas contrôlé. Un ensemble de règles de détection
traduites sous forme de requêtes s'exécute périodiquement sur ces journaux pour produire des
détections normalisées.

Un traitement planifié complète ensuite chaque détection non couverte par une règle
statique : il interroge le composant d'enrichissement sémantique, compare le résultat aux
techniques d'attaque connues, et n'inscrit un rattachement que si la similarité dépasse un
seuil déclaré --- et non codé en dur --- ce qui permet de l'ajuster sans redéploiement.

\section{Dashboard de supervision sécurité}

Le tableau de bord de supervision n'accède ni directement à l'entrepôt de détection ni au
composant d'enrichissement : il passe systématiquement par l'interface applicative
intermédiaire, qui applique un contrôle d'accès par rôle (administration, analyse, lecture
standard). Cette contrainte évite qu'une interface destinée à la consultation devienne un
chemin d'accès direct aux données sensibles.

\todo{Décrire les composants réellement livrés du tableau de bord (matrice de couverture
des techniques d'attaque, chronologie d'incident, panneau de détail, vulnérabilités
priorisées) une fois leur état d'avancement confirmé.}

\section{Planning du PFE}

\begin{attentedoc}
Le guide pédagogique demande un planning sur une page complète en fin de description du
travail réalisé. À construire à partir du calendrier réel du stage (dates de début/fin de
chaque phase), pas du calendrier indicatif de la méthodologie de cadrage --- les deux
peuvent diverger.
\end{attentedoc}

% \begin{figure}[htbp]
%   \centering
%   \includegraphics[width=\linewidth]{figures/planning_pfe.pdf}
%   \caption{Planning réel du PFE}
%   \label{fig:planning}
% \end{figure}
